% SPDX-License-Identifier: LPPL-1.3c
% Copyright (C) 2026 Christophe Jorssen
% Author and Current Maintainer: Christophe Jorssen <christophe.jorssen@gmail.com>
% LPPL maintenance status: maintained
% This file is part of luafemm. See LICENSE and LICENSES.md for its terms.
\section{Parametric components as nodes}
\label{sec:components}
A component is a genuine PGF node with a fixed physical geometry and named
anchors. One node can declare several material regions. Its text and labels
are annotations: their font, width and padding never resize its magnetic
parts. Additional E and toroidal shapes, independent section widths and ideal-mode
mean contours are described in Section~\ref{sec:ideal}.
Components use the same region, mesh, solver and cache infrastructure
as explicitly drawn paths; the chosen field model determines their computational boundary conditions.

\begin{key}{/tikz/femm/u core=\marg{options} (default \normalfont empty)}
Select a U-shaped core open towards local $+y$. The node registers one
simple closed iron contour. It has no armature, winding or gap refinement
regions. The default dimensions match the core in the opening tutorial.
\end{key}
\begin{key}{/tikz/femm/u electromagnet=\marg{options} (default \normalfont empty)}
Select a U core, an armature above its two poles, and two winding sections
around the left leg. A single node registers four material regions and,
by default, two air regions for gap refinement.
The armature has the same width as the core; both gaps have the same height.
\end{key}

Put either key on a node inside a picture with |femm/problem|. It selects
the internal component shape and enables |transform shape|. Declare every
component and additional physical path before requesting a mesh or a field.
Components do not solve during construction. They may be named or anonymous.
Ordinary PGF grouping applies; a reusable TikZ style can set their options.

\begin{codeexample}[width=6cm]
\begin{tikzpicture}[femm/problem={
  xmin=-60,xmax=60,ymin=-35,ymax=35},
  scale=.5]
  \node[femm/u core={
    width=30,height=25,
    leg thickness=5}] (A) at (-20,0) {};
  \node[femm/u core={
    width=30,height=25,
    leg thickness=5,
    core style={fill=blue!15}},
    rotate=15] (B) at (20,0) {};
  \draw[red,thick]
    (A.left pole)--(A.right pole);
  \draw[blue,thick]
    (B.left pole)--(B.right pole);
\end{tikzpicture}
\end{codeexample}

\subsection{Geometry, materials and excitation}
The following subkeys belong to |/femm/component| and are supplied inside
either component key. All lengths are numeric literals in model millimetres,
not TeX dimensions; no |mm| suffix is needed. Values are finite numbers,
not PGF expressions. Defaults are inherited within the usual TeX scope.
Electromagnet-only options have no effect on |femm/u core|.

\begin{key}{/femm/component/width=\meta{number} (initially 80)}
Outer width of the U core and armature. It must exceed twice the leg thickness.
The winding extends outside this width on the left.
\end{key}
\begin{key}{/femm/component/height=\meta{number} (initially 60)}
Outer height of the U core alone, from its underside to its pole faces.
It must exceed the leg thickness.
\end{key}
\begin{key}{/femm/component/leg thickness=\meta{positive number} (initially 15)}
Common thickness of both legs and the bottom yoke. The window has width
|width-2*leg thickness| and height |height-leg thickness|.
\end{key}
\begin{key}{/femm/component/gap=\meta{positive number} (initially 3)}
Height of each air gap between a pole and the armature; electromagnet only.
Zero is rejected in this initial component interface.
\end{key}
\begin{key}{/femm/component/armature thickness=\meta{positive number} (initially 15)}
Height of the separate armature; electromagnet only. This is an independent
parameter: changing the leg thickness does not change this default.
\end{key}
\begin{key}{/femm/component/coil width=\meta{positive number} (initially 6)}
Width of each rectangular winding section; electromagnet only.
\end{key}
\begin{key}{/femm/component/coil height=\meta{positive number} (initially 26)}
Height of each winding section; electromagnet only. Both sections must fit
between the top of the yoke and the pole faces.
\end{key}
\begin{key}{/femm/component/coil clearance=\meta{nonnegative number} (initially 2)}
Distance from each winding section to its adjacent face of the left leg.
The inner section must fit inside the U window; electromagnet only.
\end{key}
\begin{key}{/femm/component/coil center y=\meta{number} (initially 3)}
Vertical position of the two section centres relative to |core center|,
positive towards the poles. It must keep the entire winding above the yoke
and below the pole faces; electromagnet only.
\end{key}
\begin{key}{/femm/component/ampere turns=\meta{number} (initially 6000)}
Signed $NI$ in A\,turns, for the complete winding. A positive value produces
positive $J_z$ in the outer (left) section and negative $J_z$ in the inner
section. The integral of $J_z$ over each section is respectively $NI$ and
$-NI$; their net current is zero. Negative values reverse the excitation;
zero leaves a passive winding.

The package computes $J_z=NI/S$ using the area of the actual captured
polygon in square metres. Scaling or deforming the node therefore preserves
ampere-turns and adjusts current density. This convention differs from a raw
|femm/current density| on a path, which holds density fixed.
The winding is homogenised, not drawn or solved turn by turn.
\end{key}
\begin{key}{/femm/component/core material=\meta{material} (initially pure-iron)}
Material of the U core, resolved through the existing region material
selector. Catalogue identifiers and custom |femm/materials/| styles are accepted.
\end{key}
\begin{key}{/femm/component/armature material=\meta{material} (initially pure-iron)}
Material of the armature; electromagnet only.
\end{key}
\begin{key}{/femm/component/coil material=\meta{material} (initially copper)}
Material of both winding sections; electromagnet only. Excitation is determined
by |ampere turns|, not by a current density stored in the material catalogue.
The core, armature and gap regions have zero imposed current density.
\end{key}
\begin{key}{/femm/component/magnetization angle=\meta{number} (initially 0)}
Direction of any coercive material in degrees from the component's local
$+x$ axis. It applies to material parts and has no effect on passive materials.
The direction is transformed into the model frame when the node is placed;
rotating the component rotates its prescribed magnetisation. This differs
from |femm/region/magnetization angle|, which is already a model-frame angle.
\end{key}
\begin{key}{/femm/component/mesh size=\meta{nonnegative number} (initially 0)}
Local spacing for the core, armature and winding. Zero inherits the problem
spacing. As with ordinary regions, positive values refine rather than coarsen.
Mesh lengths are model-frame values; node scaling does not scale these keys.
\end{key}
\begin{key}{/femm/component/gap mesh size=\meta{nonnegative number} (initially 1)}
Local spacing in both air gaps; electromagnet only. A positive value adds
two air refinement rectangles. In open mode zero omits these two declarations and uses
the ordinary global mesh in the gaps. In ideal mode the two gaps remain
part of the confined domain, using the global mesh spacing.
\end{key}

\subsection{Anchors and placement}
The |center| anchor is the centre of the rectangular envelope of all parts,
including the exterior winding. In an electromagnet it differs from
|core center|. Use |anchor=core center| to place the U at the same coordinates
as the tutorial's explicit polygon. Text is centred at |window center|
independently of the anchor used for placement.

\begin{center}
\begin{tabular}{p{.42\linewidth}p{.49\linewidth}}
\toprule Anchor & Meaning in local component coordinates \\
\midrule
|center| & Centre of the component envelope \\
|north|, |south|, |east|, |west| & Midpoints of envelope sides \\
|north east|, |north west|,\newline |south east|, |south west|
  & Corners of the envelope \\
|core center| & Centre of the U's outer rectangle \\
|window center| & Centre of the open window; text location \\
|yoke center| & Centre of the bottom iron bar \\
|left pole|, |right pole| & Centres of the two pole faces \\
|left pole inner|, |left pole outer| & Inner/outer endpoints of the left pole \\
|right pole inner|, |right pole outer| & Inner/outer endpoints of the right pole \\
|left armature|, |right armature| & Points facing the pole centres across the gaps \\
|left gap|, |right gap| & Centres of the gaps \\
|armature center| & Centre of the armature rectangle \\
|coil positive|, |coil negative| & Centres of the outer/inner winding sections \\
\bottomrule
\end{tabular}
\end{center}
The last four rows are available only on an electromagnet; requesting one on
a bare core is an error. The coil anchor names describe the sign for positive
|ampere turns|; their geometric locations do not swap when excitation reverses.
An automatic connection such as |(M)--(N)| meets the rectangular envelope,
which may enclose air. Use physical anchors for probes or exact pole contacts.

\begin{codeexample}[code only]
\node[femm/u electromagnet={gap=3,ampere turns=6000},
  anchor=left pole,rotate=20] (M) at (0,0) {};
\draw[femm/profile={name=left-gap,samples=101},red,dashed]
  (M.left pole)--(M.left armature);
\end{codeexample}

Nodes support |at|, |anchor|, rotations and local affine transformations,
including inherited transformations because |transform shape| is enabled.
Keep that option enabled. Whole-picture transforms remain display transforms,
subject to the capture precision described in Section~\ref{sec:model}.
Changing a physical transform or component dimension invalidates the cache.
Changing text or painting styles does not. Very small rounding differences
under display transformations may conservatively cause a cache miss.

Dimensions use the original picture coordinate basis. Changing local |x|
or |y| vectors does not resize an existing shape: use component dimensions or
node transformations. |minimum width|, |minimum height|, |inner sep| and
|outer sep| do not alter its fixed physical outline or anchors.
Positioning-library operations such as |right=10mm of M| use the envelope.
Delayed placement, including matrix-managed nodes, is rejected: physical
regions must be captured at their final position, once.

\subsection{Drawing styles and a complete example}
\begin{keylist}[/femm/component]{core style=\marg{TikZ options},
  armature style=\marg{TikZ options},coil style=\marg{TikZ options},
  gap style=\marg{TikZ options}}
Append painting options to the corresponding part style within the current
scope; no extra options are appended initially. Use colours, line widths,
opacity or dash patterns here. Geometry-changing keys, path decorations
that replace contours, and physical keys inside these styles are unsupported:
use the explicit component parameters for physical changes.
\end{keylist}
\begin{stylekey}{/tikz/femm component core}
Base core style: |draw=black!70,fill=gray!15|.
\end{stylekey}
\begin{stylekey}{/tikz/femm component armature}
Base armature style: |femm component core|. Thus it inherits core-style changes
unless explicitly overridden.
\end{stylekey}
\begin{stylekey}{/tikz/femm component coil}
Base winding style: |draw=orange!70!black,fill=orange!25|.
\end{stylekey}
\begin{stylekey}{/tikz/femm component gap}
Empty initially: refinement rectangles are invisible. A |gap style| has
no effect when |gap mesh size=0| because no gap paths are emitted.
\end{stylekey}
These styles paint the individual parts. Node-level |draw| or |fill| does not
replace their styles. Ordinary node text options and |label| remain available.
Each component is registered once, in declaration order: later physical
regions may still override it according to the ordinary overlap rule.

The file \nolinkurl{examples/u-electromagnet-node.tex} builds the tutorial
electromagnet with one node, locates its measurement path with gap anchors,
and plots $|\mathbf B|$ and $B_y$. It enables the persistent cache and starts
with a coarse mesh. The core is a single concave polygon; all other regions
are separate simple polygons. Compound material paths are available
independently (Section~\ref{sec:domains}); these node templates do not need
them. Unmeshed exclusions and rotor templates remain unsupported.

The complete source below uses node anchors for the gap scan and labels.
Changing a component dimension moves these locations with the physical parts.
The result includes the field and the PGFPlots scan produced by that source.
\luafemmexample{u-electromagnet-node}
